\documentclass[12pt, a4paper]{article}
\usepackage[T2A]{fontenc}
\usepackage[utf8]{inputenc}
\usepackage[serbianc]{babel}
\usepackage{xcolor}
\usepackage{graphicx}
\usepackage{wrapfig}
\usepackage{amsmath}
\usepackage{amssymb}
\usepackage{amsthm}
\newtheorem{teorema}{Теорема}
\theoremstyle{definition}
\newtheorem{definicija}{Дефиниција}
\newtheorem{lema}{Лема}
\begin{document}
\begin{titlepage}
    \newcommand{\HRule}{\rule{\linewidth}{0.4mm}}
    \center    	
    \textsc{\LARGE Математички факултет}\\[3cm]     	
    \textsc{\Large Семинарски рад}\\[0.1cm]    	
    \textsc{\large из увода у информатику}\\[0.5cm]     	
    \HRule\\[0.4cm]   	
    {\LARGE\bfseries Водич за предмет: Дискретне структуре 1}\\[0.2cm]     	
    \HRule\\[2cm]   
    \vspace{17\baselineskip}
    		 	
    	\begin{minipage}{0.4\textwidth}
    		\begin{flushleft}
    			\large
    			\textit{Студент}\\
    			Софија Топаловић 170/2024
    		\end{flushleft}
    	\end{minipage}
    \hspace*{1cm}
    	\begin{minipage}{0.4\textwidth}
    		\begin{flushright}
    			\large
    			\textit{Професор}\\
    			Данијела Симић
    		\end{flushright}
    	\end{minipage}
    	
    	
    \vfill\vfill\vfill\vfill   	
      {\large Београд, \today}   	
    \vfill  	
    \end{titlepage}

\tableofcontents
\newpage

\section{Увод и организација курса}
У оквиру овог рада анализираћемо кључне концепте из предмета \textcolor{blue}{Дискретне структуре 1} ослањајући се искучиво на достављене наставне материјале. Према званичним информацијама о полагању, испит садржи делове који се вреднују одређеним бројем бодова.

Садржај испита обухвата следеће делове:
\begin{enumerate}
    \item \textbf{Писмени испит} -- обавезан део који се састоји од задатака обрађених на курсу и носи укупно 70 поена.
    \item \textbf{Теоријски тест} -- обавезан део који садржи 20 кратких питања (дефиниције, искази теорема и примери) и носи 20 поена.
    \item \textit{Усмени испит} -- необавезан део испита који носи преостале бодове и утиче на коначну оцену.
\end{enumerate}

\subsection{Табела бодовања}
Минимум за пролаз наставног курса захтева остваривање прагова на обавезним деловима. У табели испод је приказана тачна расподела бодова према званичном документу:

\begin{table}[h]
    \centering
    \begin{tabular}{|l|c|c|}
        \hline
        \textbf{Компонента испита} & \textbf{Максимално поена} & \textbf{Праг пролазности} \\
        \hline
        Писмени испит & 70 & 30 \\
        Теоријски тест & 20 & 10 \\
        Бонус бодови & 10 & / \\
        \hline
    \end{tabular}
    \caption{Структура поена на Дискретним структурама 1}
\end{table}

\section{Исказна логика и задовољивост}
Формални математички системи почињу дефинисањем формула. Исказне формуле се граде од исказних слова, логичке константе контрадикције ($\perp$) и импликације ($\rightarrow$).

\begin{wrapfigure}{l}{0.4\textwidth}
    \centering
    \includegraphics[width=0.35\textwidth]{slike/slika.png}
    \caption{Симболични приказ задовољивости}
\end{wrapfigure}
Појам валуације омогућава одређивање тачности формула. Кажемо да је скуп формула задовољив ако постоји валуација при којој су све формуле тог скупа тачне. Један од најзначајнијих теоријских резултата у овој области је \textcolor{teal}{Теорема компактности} која повезује цео скуп са његовим коначним деловима, омогућавајући анализу бесконачних конфигурација кроз коначне подскупове.

\section{Теоријски модели и формална тврђења}
Дискретна математика користи прецизно дефинисана окружења за исказивање тврђења, лема и дефиниција.

\subsection{Дефиниције и функционалне релације}
Пратећи материјале о бинарним релацијама и функцијама, уводимо строгу дефиницију функционалности.

\begin{definicija}
Бинарна релација $f \subseteq A \times B$ је функционална или парцијална функција ако важи следећи исказ:
\begin{equation}
    (\forall a \in A)(\forall b_1, b_2 \in B)(a f b_1 \wedge a f b_2 \rightarrow b_1 = b_2)
\end{equation}
\end{definicija}

Ова дефиниција осигурава да из сваког елемента домена излази највише једна стрелица ка кодомену.

\subsection{Леме и теореме из елементарне аритметике}
Релаcija дељивости на скупу целих бројева $\mathbb{Z}$ заснива се на постојању количника. За запис особина користимо формална математичка тврђења.

\begin{lema}
Релација дељивости поседује особину рефлексивности, што значи да за сваки цео број важи исказ $n|n$, као и особину транзитивности.
\end{lema}
Ове особине можемо прегледно приказати:
\begin{itemize}
    \item \textbf{Рефлексивност:} за сваки цео број $n$ важи $n|n$.
    \item \textbf{Транзитивност:} ако $a|b$ и $b|c$, онда важи $a|c$.
\end{itemize}
Поред релација, елементарна аритметика се бави и теоријом конгруенција и простих бројева, где се истиче Вилсонова теорема.

\begin{teorema}
Нека је $n$ природан број већи од 1. Број $n$ је прост ако и само ако важи следећа конгруенција по модулу $n$:
\begin{equation}
    (n-1)! \equiv_n -1
\end{equation}
\end{teorema}

\section{Закључак}
У овом раду су успешно приказани кључни теоријски концепти из наставних материјала за предмет Дискретне структуре 1. Коришћењем система \LaTeX\ креиран је прегледан академски документ на српском језику уз поштовање свих правила ћириличне типографије. Кроз рад су имплементиране сложене математичке формуле, табеларни и листични прикази, као и напредна окружења за дефиниције, леме и теореме.

\end{document}